% Spoken duality scripts, based on the standard-duality recording.
% Finite-volume/localization notation is suppressed in the talk.

\newcommand{\DualityProcessScript}{%
{\color{talkblue}\bfseries Slide 10: Duality: tracing dependence backwards}\par
{\small 8:00--9:25 (85 seconds)}\par\medskip
Our proof starts from duality because it is a standard way to study convergence without coupling and without a priori knowing the limiting measure. Classically, this proves ergodicity when interactions are sufficiently weak compared with independent updates.

\keyscript{The dual consists of a finite set $A_s$ and a sign $\sigma_s$.} Starting from $A$, think of $A_s$ as tracing backwards which earlier spins can influence the spins in $A$ at time t.

The three types of updates are deaths, splits and births. The picture shows an example trajectory starting from sites three and four. A death removes site four. A split removes site three and adds sites two and four. A birth then adds sites one and three, while leaving sites two and four active.

\keyscript{Duality gives the Feynman--Kac representation below.} $\chi_A$ is the product of the spins in $A$. The representation also features a sign process $\sigma_t$ and an exponential weight involving a potential $V$.%
}

\newcommand{\DualityClassicalScript}{%
{\color{talkblue}\bfseries Slide 11: The usual duality argument}\par
{\small 9:25--10:25 (60 seconds)}\par\medskip
For ergodicity, we want $P_t\chi_A$ to converge to a constant with respect to $\eta$, for every finite $A$.

\keyscript{Duality lets us split $P_t\chi_A$ into a term which does not depend on the initial configuration and a term which is negligible if the process is ergodic.} We make this split according to whether $A_t$ is empty.

The classical argument assumes that the potential is bounded above by a strictly negative constant while the dual set is nonempty. On that event, the exponential weight is at most $e^{-\gamma t}$. The second term vanishes, giving a limit independent of the initial configuration.%
}

\newcommand{\DualityObstructionScript}{%
{\color{talkblue}\bfseries Slide 12: Why the standard estimate fails}\par
{\small 10:25--11:00 (35 seconds)}\par\medskip
For weakly noisy FA-1f, this estimate fails. \keyscript{The potential is positive rather than negative, so the exponential weight can be very large.}

\keyscript{To show that the term depending on $\eta$ is negligible, we have to control the cancellations from $\sigma_t$.} The large weight makes that control particularly demanding because even rare events can make substantial contributions.

Next we consider how to average over part of the graphical construction to capture these cancellations.%
}
